\section{E. The special torus group, arithmetic braid action, and surface germs}
\label{sec:E}

All mapping classes in this section preserve the specified puncture.
The superscript $+$ means preservation of the orientation.  A tangent
\emph{ray}, its universal cover, a nonzero tangent vector, and an outer
peripheral conjugacy class are different data.  We state each group
using the data it actually preserves.  In particular, $\widehat{\mathbb Z}$
and the discrete group $\mathbb Z$ are not interchangeable in the
covering constructions below.

\subsection{The real and integral covering groups}

\begin{statement}{E.1}{Proved}{The lifted circle action and the signed kernel}
Put $G=\{u\in\mathrm{GL}_2(\mathbb R):\det u=\pm1\}$, and identify
$S^1$ with $\mathbb R/\mathbb Z$ by $t\mapsto e^{2\pi it}$.  Define
\[
 \widetilde G=\{(u,h):u\in G,\ h:\mathbb R\xrightarrow{\sim}\mathbb R,
 \ e^{2\pi i h(t)}=u(e^{2\pi it})/|u(e^{2\pi it})|\},
\]
with product $(u,h)(v,k)=(uv,h\circ k)$ and the covering topology.
There is an exact sequence
\[
 1\longrightarrow\mathbb Z\langle d\rangle\longrightarrow
 \widetilde G\longrightarrow G\longrightarrow1,
 \qquad d(t)=t+1,
\]
and $(u,h)d(u,h)^{-1}=d^{\det u}$.  The inverse image of
$\mathrm{SL}_2(\mathbb R)$ is its universal covering group and is
homeomorphic to $\mathbb R^3$.  The inverse image
$\widetilde{\mathrm{GL}_2(\mathbb Z)}$ is a discrete group.  Its
extension by $\mathbb Z$ is central only over the positive subgroup.
\end{statement}
\begin{proof}
The normalized linear action on $S^1$ has degree $\det u\in\{\pm1\}$.
Every lift satisfies $h(t+1)=h(t)+\det u$, and any two lifts differ
by an integer translation.  These facts prove the group law, exactness,
and conjugation formula.  The polar decomposition retracts
$\mathrm{SL}_2(\mathbb R)$ onto $\mathrm{SO}(2)$, whose standard
rotation loop has winding number one on $S^1$.  The indicated covering
therefore kills exactly the whole fundamental group.  More explicitly,
the Iwasawa coordinates give
$\mathrm{SL}_2(\mathbb R)\simeq S^1\times\mathbb R^2$ as a space;
lifting the first coordinate gives $\mathbb R^3$.
The integral matrices form a discrete subgroup downstairs and a
covering is locally a homeomorphism, proving the last topological
assertion.  A reflection acts on the kernel by inversion, so that
kernel is not central in the full group.  This is the precise
construction of \S53.2--3, C4 scans 228--231, fols.~761'--763.
\end{proof}

\begin{statement}{E.2}{Proved}{The positive special group is $B_3$, with centre-square kernel}
Let
\[
 A=\begin{pmatrix}1&1\\0&1\end{pmatrix},\quad
 B=\begin{pmatrix}1&0\\-1&1\end{pmatrix}.
\]
The inverse image of $\mathrm{SL}_2(\mathbb Z)$ in E.1 is
\[
 B_3=\langle a,b\mid aba=bab\rangle
     =\langle p,q\mid p^3=q^2\rangle,
 \quad p=ab,\quad q=aba.
\]
Its centre is the infinite cyclic group generated by
$z=(ab)^3=p^3=q^2$.  Under $a\mapsto A$, $b\mapsto B$,
\[
 B_3/\langle z\rangle=\mathrm{PSL}_2(\mathbb Z),\qquad
 B_3/\langle z^2\rangle=\mathrm{SL}_2(\mathbb Z).
\]
Consequently the kernel of the covering of $\mathrm{SL}_2(\mathbb Z)$
is $\langle z^2\rangle$, not its whole centre.  With the clockwise
matrix conventions displayed here, $z^2=d^{-1}$ in E.1.
\end{statement}
\begin{proof}
The Tietze change of generators has inverse
$a=p^{-1}q$, $b=q^{-1}p^2$; substitution proves the equivalence of
presentations.  The common power $p^3=q^2=z$ is central.  Quotienting
by it gives $C_3*C_2$.  This is $\mathrm{PSL}_2(\mathbb Z)$: the
usual fundamental domain $|\Re w|\leq1/2$, $|w|\geq1$ in the upper
half-plane has the two elliptic vertex stabilizers of orders $2,3$,
and its side identifications give their free product.  Equivalently,
the dual Farey tree has quotient a single edge, with these two vertex
stabilizers and trivial edge stabilizer.  The normal form theorem for
a free product also proves its centre is trivial.  Thus every central
element upstairs belongs to $\langle z\rangle$.  The homomorphism
$a,b\mapsto1$ to $\mathbb Z$ sends $z$ to $6$, so $z$ has infinite
order and this is exactly the centre.

The matrices $P=AB$ and $Q=ABA$ have
$P^3=Q^2=-I$ and generate $\mathrm{SL}_2(\mathbb Z)$, by elementary
row reduction.  The resulting surjection from $B_3$ is an isomorphism
modulo $\langle z\rangle$ onto the projective group.  Its kernel is
therefore contained in $\langle z\rangle$ and is precisely
$\langle z^2\rangle$.

For the covering-group identification, lift $P$ with rotation number
$-1/6$ and $Q$ with rotation number $-1/4$.  Their third and second
powers are the same lift of $-I$, namely $t\mapsto t-1/2$.
They generate the full inverse image of the modular group because
they generate downstairs and their sixth/fourth powers generate the
deck kernel.  The induced map from $B_3$ has kernel contained in
$\langle z^2\rangle$; it is injective there, since $z^2$ is the
nontrivial deck translation $d^{-1}$.  This proves the identification
without assuming the braid presentation from the manuscript.
In its generators $\rho=p^{-1}$, $\sigma=q$, the relation is
$\rho^3\sigma^2=1$, as at scan 234, equation (25).
\end{proof}

\begin{statement}{E.3}{Proved}{The full special torus group and its reflection presentation}
Let $\operatorname{Diff}(T^2,0)$ carry the smooth topology.  Lift its
action on the circle of tangent rays at $0$ to the universal cover
of that circle, and define $ST_{1,1}$ to be the component group of
the resulting group of lifts.  Then
\[
 ST_{1,1}\cong\widetilde{\mathrm{GL}_2(\mathbb Z)}
 \cong B_3\rtimes\langle r\rangle,
\]
where
\[
 r^2=1,\qquad rar=b^{-1},\qquad rbr=a^{-1}.
\]
Thus $ST_{1,1}^{+}=B_3$, proving Lead~2 with its exact central
normalization.  In the generators of E.2 the full presentation is
\[
 \langle\rho,\sigma,r\mid
 \rho^3\sigma^2=1,\ r^2=1,
 r\rho r=\rho^{-1},\ r\sigma r=\sigma^{-1}\rangle.
\]
Equivalently, with $r_0=r$, $r_1=\rho^{-1}r$, $r_\infty=\sigma r$,
\[
 \langle r_0,r_1,r_\infty\mid r_0^2=r_1^2=r_\infty^2=1,
 (r_0r_1)^3(r_\infty r_0)^2=1\rangle.
\]
Conjugation by any negative class inverts $z$ and $z^2$.
\end{statement}
\begin{proof}
The component group of $\operatorname{Diff}(T^2,0)$ is
$\mathrm{GL}_2(\mathbb Z)$: the action on $H_1$ is surjective by
linear torus maps, and a marked torus diffeomorphism acting trivially
on $H_1$ is isotopic to the identity relative to the marking.  The
identity component is contractible by
\citet[Theorem~1(b)]{EarleEells1967}.  The exact homotopy sequence of
the pulled-back circle covering consequently gives
\[
 1\to\mathbb Z\to ST_{1,1}\to\mathrm{GL}_2(\mathbb Z)\to1.
\]
The linear representatives identify this extension with that in E.1:
there is one linear representative in each component downstairs,
and no loop in the identity component can identify different lifts.

Take the linear reflection $R(x,y)=(y,x)$ and its order-two lift
$r(t)=1/4-t$.  Choose $a$ as the lift of $A$ fixing $t=0$, and $b$
as the lift of $B$ fixing $t=1/4$.  The equalities
$RAR=B^{-1}$, $RBR=A^{-1}$, together with these fixed points, give
the stated lift relations.  They define the split extension of
E.2 by $C_2$, so give the whole presentation.  Substituting
$\rho=(ab)^{-1}$, $\sigma=aba$ proves the second presentation;
substituting the three involutions proves the third.  Finally
$r(ab)^3r=(b^{-1}a^{-1})^3=z^{-1}$.
The positive group can also be described as the mapping class group
of a torus with one boundary circle fixed pointwise.  An orientation
reversing map cannot fix that boundary pointwise; the full signed
statement above uses tangent lifts, not that boundary convention.
This proves scans 234--239, equations (25), (28), and (30).
\end{proof}

\begin{statement}{E.4}{Refuted}{Three literal formulas need correction}
With the conventions visibly used in \S53, the following formulas
cannot stand literally: the normalization $u(z)/|z|$ on scan 228;
the simultaneous shift of lifts by exponents $(2n,3n)$ in scan 232,
equation (18); and the formula containing $2i\pi(1/8-t)$ for a map
$\mathbb R\to\mathbb R$ on scan 233, equation (23).
Their corrected forms are, respectively,
\[
 u(z)/|u(z)|,\qquad
 (\rho d^{2n},\sigma d^{-3n}),\qquad
 \widetilde{(x,y)\mapsto(y,x)}(t)=1/4-t.
\]
The last lift fixes $t=1/8$, the ray of angle $\pi/4$.
\end{statement}
\begin{proof}
For $u=\operatorname{diag}(2,1/2)$ and $z=1$, the first proposed
expression has modulus $2$ instead of $1$.  The corrected normalized
map has modulus one.  If $\rho^3\sigma^2=1$, replacing the lifts by
$\rho d^a,\sigma d^b$ changes this product by $d^{3a+2b}$.
Its full solution lattice is $(a,b)=(2n,-3n)$; the displayed positive
sign would produce $d^{12n}$.  Finally swapping real and imaginary
parts changes angle $2\pi t$ to $\pi/2-2\pi t$.  Its lift is therefore
$1/4-t$, with fixed point $1/8$; no factor $2i\pi$ can occur in its
real-valued formula.  The problematic signs and coefficients on
scans 232--233 were checked against the actual images, not silently
repaired in the transcription.  Scan 231's kernel generator must
likewise mean the full-turn lift of $I$, not the half-turn lift of
$-I$; E.2 explains the distinction independently of a doubtful prime.
\end{proof}

\subsection{The question mark in the early draft}

\begin{statement}{E.5}{Proved}{The early draft gives the universal-cover extension with its negative kernel convention}
Interpret the diagram of cote 140-1, scan 29, together with the
kernel convention recorded on scans 35--36, now inspected in the
original archive images \citep{Grothendieck1981Draft}, and the
embedding specified on scans 20 and 36:
\[
 \mathrm{SL}_2(\mathbb Z)\hookrightarrow
 \mathrm{PSL}_2(\mathbb Z)\times\mathbb Z/4,
 \qquad s\longmapsto(\bar s,\operatorname{ab}(s)\bmod4).
\]
The group denoted by ``?'' is $B_3$ as an extension of the
\emph{specified} $\mathrm{SL}_2(\mathbb Z)$ by the
\emph{specified} $\mathbb Z$: its kernel is the square of the
braid centre, and it is the universal-cover extension of E.2.
The sign is essential.  The draft uses the kernel generator
$\omega'=\omega^{-1}$ for the extension over
$\mathrm{PSL}_2(\mathbb Z)$; pushing out the opposite generator
would instead give class $5$ in
$H^2(\mathrm{SL}_2(\mathbb Z),\mathbb Z)\cong\mathbb Z/12$,
although its middle group would still be abstractly $B_3$.
Here E.2's specified extension is class $1$.
\end{statement}
\begin{proof}
Start from $H=\langle p,q\mid p^3=q^2=\omega\rangle$, with
$\operatorname{ab}(p)=2$, $\operatorname{ab}(q)=3$.  Scan 36's
image and checked transcription explicitly give the pushout presentation
\[
 H_4^-=\langle p,q,t\mid t\text{ central},\ p^3=q^2=t^{-4}\rangle,
 \qquad t^4=\omega'.
\]
Map $H_4^-$ onto $\mathrm{PSL}_2(\mathbb Z)\times C_4$ by
$p\mapsto(\bar p,0)$, $q\mapsto(\bar q,0)$, $t\mapsto(1,1)$.
The kernel is $\langle t^4\rangle$.  Its pullback to the displayed
copy of $\mathrm{SL}_2(\mathbb Z)$ is the index-two subgroup
\[
 P^-=\langle p,b,v\mid v\text{ central},\ p^3=v^{-2},\ b^2=v^{-1}\rangle,
 \qquad b=qt,\quad v=t^2.
\]
This presentation follows either by the two-coset Schreier rewriting
with representatives $1,t$, or by projecting to the two parity
characters.  Set $\alpha=pv$, $\beta=bv=qt^3$.  Then
$\alpha^3=\beta^2=v$, $p=\alpha^{-2}$ and $b=\beta^{-1}$.
These inverse substitutions give
$P^-=\langle\alpha,\beta\mid\alpha^3=\beta^2\rangle\cong B_3$.
The quotient sends $\alpha$ and $\beta$ to the specified modular
generators $R$ and $S$, since their two coordinates are
$(\bar p,2)$ and $(\bar q,3)$.  Its specified kernel is
$\langle v^2\rangle=\langle t^4\rangle$, exactly the centre-square
kernel of E.2.  The maps $p\mapsto pt^2$, $q\mapsto qt^3$,
$\omega\mapsto t^2$ and $\omega^2\mapsto t^4=\omega'$ are also
the maps displayed in the transcription of scan 36.

For comparison, if one ignores $\omega'=\omega^{-1}$ and replaces
$t^{-4}$ by $t^4$, the resulting pullback has presentation
\[
 P^+=\langle p,b,v\mid v\text{ central},\ p^3=v^2,\ b^2=v^3\rangle.
\]
Now $p^{-1}v$ and $bv^{-1}$ again have common cube and square $v$,
so the middle group is abstractly $B_3$.  But the specified lifts
of $R,S$ are
$\widetilde R=pv$, $\widetilde S=bv$.  They satisfy
$\widetilde R^3=\widetilde S^2=v^5$, and hence
$\widetilde R^6=\widetilde S^4=(v^2)^5$.
The amalgam $\mathrm{SL}_2(\mathbb Z)=C_6*_{C_2}C_4$ gives
\[
 H^2(\mathrm{SL}_2(\mathbb Z),\mathbb Z)
 =\{(u\bmod6,w\bmod4):u\equiv w\pmod2\}.
\]
Indeed the group-cohomology exact sequence for the amalgam is
injective in degree two because $H^1(C_2,\mathbb Z)=0$, and its image
is the kernel of the two restrictions to $H^2(C_2,\mathbb Z)$.
One can obtain the same description directly by choosing lifts of
the two cyclic generators and comparing their common central
involution.  The Chinese remainder identification is $\mathbb Z/12$.
The wrong-sign lift-power calculation gives
$(5\bmod6,5\bmod4)$, whereas $P^-$ and E.2 give $(1,1)$.
This explains why an abstract identification of the middle group
alone cannot decide the diagram.

\emph{Reading qualification.}  Besides the supplied images of scans
20, 27 and 29, we retrieved the original archive PDF and inspected
its pages 35--36 (bearing manuscript labels 34--35).
The visible negative exponent in scan 36's upper-right presentation
and its four displayed maps are mutually consistent.  Scan 35
visibly records $(\alpha=-1\ ?)$ and $(3-2)\eta_{12}$, supporting
class $1$.  Some overwritten prime and inverse labels remain
editorially uncertain; the proof does not need to resolve those
individual strokes because both the presentation and the maps fix
the convention.  The class-$5$ calculation records exactly what
would happen under the opposite convention, not an alternative
assertion about the intended result of the later draft.
\end{proof}

\subsection{The based normalizer and the arithmetic section}

\begin{statement}{E.6}{Known}{The standard Grothendieck--Teichmuller braid action}
For $(\lambda,f)\in\widehat{\mathrm{GT}}$, with
$f\in[\widehat F_2,\widehat F_2]$ and the usual Drinfeld
normalization, the action on $\widehat B_3$ is
\begin{equation}\label{eq:E-standard-action}
 a\longmapsto a^\lambda,\qquad
 b\longmapsto f(a^2,b^2)^{-1}b^\lambda f(a^2,b^2).
\end{equation}
It is an automorphism, sends $z=(ab)^3$ to $z^\lambda$, and
preserves the subgroup $\langle z^2\rangle\cong\widehat{\mathbb Z}$.
In the category of profinite parenthesized braids this is the
restriction of the defining $\widehat{\mathrm{GT}}$ action;
the hexagons give the braid relations and the pentagon gives the
arity-four coherence.  The formula is explicitly recorded in
\citet[(5.4.2)]{Nakamura2013Eisenstein}; the equivalent operadic
formulation, including its finite quotient conditions, is given in
\citet[\S2 and Appendix A]{DolgushevLeLorenz2024}.
\end{statement}
\begin{proof}[Explanation of the restriction and normalization]
The elementary crossing $a$ goes to its $\lambda$-power.  To express
the second crossing with the same parenthesization one changes
parentheses, whose image is the associator $f$, applies the crossing,
and changes back.  This gives the conjugation in
\eqref{eq:E-standard-action}.  The defining coherence identities
are exactly what make these assignments a functor; applying the
inverse element of $\widehat{\mathrm{GT}}$ supplies the inverse
functor, so checking a single braid relation is not being substituted
for invertibility.  The centre of $\widehat B_3$ is
$\overline{\langle z\rangle}$: after dividing by it the group is the
profinite free product $C_3\mathbin{\widehat *}C_2$, which has trivial
centre.  The procyclic centre has its full profinite topology because
the abelianization sends $z$ to $6$.  An automorphism sends $z$ to
$z^\mu$ for some $\mu\in\widehat{\mathbb Z}^{\times}$.  Abelianizing
\eqref{eq:E-standard-action} gives $6\mu=6\lambda$; multiplication
by $6$ on $\widehat{\mathbb Z}$ is injective, proving $\mu=\lambda$.
This also identifies its action on the centre-square kernel.
\end{proof}

\begin{statement}{E.7}{Proved}{The precise tangential normalizer and its quotient}
Let $F=\widehat{\langle x,y\rangle}$, $c=[x,y]$, $L=\overline{\langle c\rangle}$,
and put
\[
 A_L=\{\varphi\in\operatorname{Aut}(F):\varphi(c)=c^\mu
                 \text{ for some }\mu\in\widehat{\mathbb Z}^{\times}\}.
\]
Embed $\widehat B_3$ by the completed geometric action
$A(x)=x,A(y)=yx$, $B(x)=xy^{-1},B(y)=y$ of D.6, and define
\[
 SN_{1,1}=N_{A_L}(\widehat B_3).
\]
Then
\[
 1\to\widehat B_3\to SN_{1,1}\to\widehat{\mathrm{GT}}\to1
\]
is exact.  The standard braid action supplies a splitting, with
conjugation action E.6.  If ``loop preserving'' is required to fix
$c$ itself, the quotient is instead
$\widehat{\mathrm{GT}}_1=\ker(\lambda:\widehat{\mathrm{GT}}
\to\widehat{\mathbb Z}^{\times})$.  Thus the full statement requires
preservation of $L$, allowing its cyclotomic power automorphisms.
\end{statement}
\begin{proof}
We first justify the peripheral kernel in this passage from outer to
based automorphisms.  By \citet[\S1.3, Lemma~1]{Asada2001},
$C_F(c)=L$.  There is also a short concrete reduction in this case.
In the index-two subgroup with $x\mapsto1,y\mapsto0$ in $C_2$,
a Schreier basis is $x^2,y,xyx^{-1}$, so $c=(xyx^{-1})y^{-1}$ is
primitive and generates a procyclic free factor.  The free-factor
malnormality theorem, explicitly stated in
\citet[Proposition~4.4]{HaranJardenPop2002}, gives its centralizer
inside this subgroup.  Conjugation by $x$ negates the nonzero class
$[c]$ in its abelianization; conjugation by an element of that
subgroup acts trivially there.  No element outside it can centralize
$c$.  Finally an inner automorphism sending $c$ to $c^\mu$ must
have $\mu=1$: map to the finite Heisenberg groups over $\mathbb Z/n$,
where $c$ is central of order $n$, for every $n$.  Hence
$N_F(L)=C_F(c)=L$.

Let $O=\operatorname{Out}_{\mathrm{lac}}(F)$ preserve the peripheral
conjugacy class up to unit power, and let
$\Sigma\cong\widehat{\mathrm{SL}_2(\mathbb Z)}$ be its geometric
subgroup.  Its faithful embedding is the genus-one congruence theorem
\citep[\S7, Theorem~5]{Asada2001}.  The map $A_L\to O$ is surjective
by changing an inner representative, and its kernel is
$\operatorname{Int}(L)$.  The inverse image of $\Sigma$ is
$\widehat B_3$: E.2 and D.7 identify the kernel generator $z^2$ with
$\operatorname{Int}(c^{-1})$, and the surjection to $\Sigma$ supplies
all other lifts.  Faithfulness follows from the faithful outer
representation and the faithful action of this procyclic kernel.
The elementary inverse-image normalizer identity now gives
\[
 SN_{1,1}/\widehat B_3\cong N_O(\Sigma)/\Sigma.\tag{E.7.1}
\]

For completeness we check the substantive input to LM3 E.10.
The inverse image of $\Sigma$ in $\operatorname{Aut}(F)$ is
$E=\widehat{B_4/\langle\zeta_4\rangle}$, where
$\zeta_j=(\sigma_1\cdots\sigma_{j-1})^j$.  The Ferrari map
$\sigma_1,\sigma_2,\sigma_3\mapsto a,b,a$ sends
$\zeta_4$ to $z^2$.  Its quotient gives
\[
 1\to F\to E\to\widehat{\mathrm{SL}_2(\mathbb Z)}\to1.
\]
Here the free Ferrari kernel and its two generators are given
explicitly in \citet[(5.4.3)--(5.4.4)]{Nakamura2013Eisenstein};
the central quotient is the once-more-pointed elliptic modular
group described in his \S5.7.  Its conjugation map to
$\operatorname{Aut}(F)$ is injective: any kernel element projects
trivially to the faithful outer image $\Sigma$, hence lies in
the centre-free kernel $F$, where its conjugation is faithful.
It is onto the inverse image of $\Sigma$ because both have
kernel $\operatorname{Int}(F)$ and quotient $\Sigma$.
The kernel $K$ of the further map to
$\widehat{\mathrm{PSL}_2(\mathbb Z)}$ is
$C_2\mathbin{\widehat *}C_2\mathbin{\widehat *}C_2$ and is
characteristic in $E$ by
\citet[Proposition~3.2(iii), (3.5)]{MinamideNakamura2022}.
The subgroup $F$ is the kernel of the unique map $K\to C_2$ taking
each of the three conjugacy classes of involutions nontrivially.
The free-product torsion theorem used there makes this characterization
invariant, so $F$ is characteristic in $E$.

Conjugation by an element normalizing $\Sigma$ induces an outer
automorphism of $E$.  Its kernel modulo $\Sigma$ is trivial: if this
automorphism is inner, subtract that inner automorphism and obtain an
automorphism of $F$ commuting with every inner automorphism of $F$.
It fixes every element because $F$ has trivial centre.  Thus (E.7.1)
embeds in $\operatorname{Out}(E)$, which is
$\widehat{\mathrm{GT}}$ by
\citet[Theorem~B = Theorem~4.3 and Corollary~C]{MinamideNakamura2022}.
Conversely the standard braid action preserves the Ferrari kernel:
after applying the Ferrari map, its third-generator conjugator becomes
$f(z,a^2)=1$, since $z$ is central and $f$ is a commutator word.
It preserves the subgroup generated by $\sigma_1,\sigma_2$ and sends
the peripheral word $\zeta_3^2\zeta_4^{-1}$ to its $\lambda$-power.
Restricting to $F$ therefore lies in $SN_{1,1}$.  Its outer action on
$E$ is the specified isomorphism of the cited theorem, so this
restriction gives both surjectivity and a splitting of (E.7.1).
Restricting the multiplier to $1$ proves the final qualification.
\end{proof}

\begin{statement}{E.8}{Known}{The Tate--Weierstrass section has a Kummer correction}
Over $\mathbb Q$, the moduli scheme of elliptic curves with a nonzero
invariant differential is
\[
 M_{1,1}^{\omega}=\operatorname{Spec}
 \mathbb Q[g_2,g_3,(g_2^3-27g_3^2)^{-1}].
\]
Its geometric fundamental group is $\widehat B_3$.  The Tate curve
at the cusp, with its specified differential and Weierstrass tangent,
provides a section of the arithmetic fundamental-group sequence
\[
 1\to\widehat B_3\to\pi_1(M_{1,1}^{\omega},\bar q)
       \to G_{\mathbb Q}\to1.
\]
It embeds this arithmetic extension in $SN_{1,1}$ of E.7.  Write
$\chi$ for the cyclotomic character, $f_\sigma$ for the standard
Drinfeld associator, and $\rho_2$ for the Kummer cocycle specified by
$\sigma(2^{1/n})/2^{1/n}=\zeta_n^{\rho_2(\sigma)}$ with compatible
choices.  At the Tate--Weierstrass base point the action is
\begin{align*}
 a&\longmapsto a^{\chi(\sigma)},\\
 b&\longmapsto a^{-8\rho_2(\sigma)}
 f_\sigma(a^2,b^2)^{-1}b^{\chi(\sigma)}f_\sigma(a^2,b^2)
 a^{8\rho_2(\sigma)}.
\end{align*}
Thus applying E.6 without a path correction gives the standard braid
base point, not these specified Tate generators.
\end{statement}
\begin{proof}[Construction and comparison of the sections]
The cubic $4X^3-g_2X-g_3$ has three distinct unordered roots of sum
zero.  Its discriminant complement is the configuration space of
three unordered complex points modulo their mean translation.
The latter translation factor is contractible, so its fundamental
group is $B_3$.  The algebraic/topological comparison gives the
stated profinite geometric group.  The arithmetic construction is
\citet[\S\S2.3--2.4, 5.2, and Proposition~5.7.3]{Nakamura2013Eisenstein}:
the tangent coordinate at the identity of the universal curve supplies
the Weierstrass section of the Ferrari sequence, and the Tate
parameter supplies the section over $G_{\mathbb Q}$.  Proposition~5.7.3
identifies its image with $\langle\tau_1,\tau_2\rangle\rtimes s_1(G_{\mathbb Q})$
and fixes the two named geometric generators.

To align the free basis with D.6, one transports Nakamura's Ferrari
basis by the half twist $q=ABA$.  In his original basis the actions
of $\tau_1,\tau_2$ are $B,A$, respectively
\citep[Proposition~5.7.3]{Nakamura2013Eisenstein}; conjugation by
$q$ exchanges $A$ and $B$ and fixes $c$.  After this explicit
identification we write $a=\tau_1$, $b=\tau_2$.  The braid formulas
are unchanged, while the based free-group identification is fixed.

The path from the standard braid base point to the Tate--Weierstrass
base point is computed in that paper's (5.7.1).  Its cocycle is
$\omega_2^{4\rho_2}\omega_3^{-4\rho_2}\omega_4^{\rho_2}$,
where $\omega_2=a^2$, $\omega_3=z$, and $\omega_4=\zeta_4$.
Conjugating the standard action by its inverse gives (5.7.2).
On $\widehat B_3$, the $z$-factor is central and the $\zeta_4$-factor
acts trivially; the remaining factor is $a^{-8\rho_2}$, giving the
formula displayed.  On the Ferrari kernel $F$, where the action is
based, the full correction inside $\widehat B_3$ is
$a^{-8\rho_2}z^{4\rho_2}$.  It is a cocycle because $a$ and $z$
commute and both are raised to $\chi(\sigma)$ by the standard action.
This explains the embedding in E.7 and records a correction which
would be invisible after passing to the outer braid action.  The
faithfulness used here is the standard embedding
$G_{\mathbb Q}\hookrightarrow\widehat{\mathrm{GT}}$
recalled in \citet[\S5.1]{Nakamura2013Eisenstein}: together with
the faithful geometric kernel of E.7 it makes the whole
arithmetic extension injective.  The
arithmetic subgroup has quotient $G_{\mathbb Q}$; identifying it with
all of $SN_{1,1}$ would require surjectivity of
$G_{\mathbb Q}\to\widehat{\mathrm{GT}}$, which is not asserted.
\end{proof}

\subsection{Finite calculations with a stated exhaustive range}

\begin{statement}{E.9}{Proved}{Genuine shadows in the abelian-pure family, and the separate matrix audit}
Let $Q_j(m)=B_j/([PB_j,PB_j]PB_j^m)$ for $j\leq4$.
Every charming shadow for the compatible arity-four quotient $Q_4(m)$
is genuine, and all are
\[
 (k\bmod m,\bar1),\qquad 0\leq k<m,\quad\gcd(2k+1,m)=1.
\]
For the complete range $|Q_3(m)|\leq1000$ the executed list is
\[
\begin{array}{c|r|l|r}
 m&|Q_3(m)|&\lambda\bmod2m&\text{number of shadows}\\\hline
 1&6&1&1\\
 2&48&1,3&2\\
 3&162&1,5&2\\
 4&384&1,3,5,7&4\\
 5&750&1,3,7,9&4
\end{array}
\]
There are no other $m$ in this family under the bound, since
$|Q_3(6)|=1296$.  The program checks both hexagons, the pentagon,
full-twist powers, and the complete quotient automorphisms.
It also checks all $304192$ power automorphisms of the cyclic
quotients $C_n$, $1\leq n\leq1000$.
\end{statement}
\begin{proof}
For $N=[PB_4,PB_4]PB_4^m$, all pure linking numbers are recorded
modulo $m$.  Insertion and cabling of strands preserve this condition,
and insertion of three strands detects each of their linking numbers;
thus the induced arity-three quotient is exactly $Q_3(m)$, and
$N_{\mathrm{ord}}=m$.  In the abelian pure quotient a commutator
associator has image $1$.  Both hexagons reduce to the conjugation
identities between the three pure generators, and the pentagon is
literally $1=1$.  The action on the pure abelian subgroup is scalar
multiplication by $2k+1$, so its surjectivity is exactly the displayed
coprimality condition.  This is also
\citet[Theorem~B.2]{DolgushevLeLorenz2024}.  To check genuineness
without an inference from finite relations, choose an actual Galois
element whose cyclotomic character has the prescribed odd unit
residue modulo $2m$; cyclotomic surjectivity supplies it.  Its
commutator associator is trivial in this quotient and gives the
required shadow.

The representation used in \texttt{e\_genuine\_shadows.py} records an
endpoint permutation and a signed crossing vector modulo $2m$.
Each pair's parity is prescribed by the permutation, so
$|Q_j(m)|=j!m^{\binom j2}$.  Multiplication is the sum of crossing
vectors after permutation of strand labels.  Consequently
$(p,v)\mapsto(p,\lambda v)$ is a homomorphism, is bijective for the
listed units, and commutes with cabling, which duplicates crossing
coordinates.  The program enumerates all elements of the indicated
$Q_3(m)$ and verifies every Cayley edge against the substituted
generators; it independently evaluates both hexagon words.
The arity-four pentagon has trivial associator on every side.
The cyclic computations in \texttt{e\_finite\_braids.py} are exhaustive
for $n\leq1000$, and their genuineness follows in the same way from
the cyclotomic character.  These are rigorous finite checks with
specified quotient families, not an enumeration of every finite
group of order at most $1000$.
\end{proof}

A second, deliberately separate computation enumerates matrices in
$\mathrm{SL}_2(\mathbb Z/n)$, $2\leq n\leq10$, takes
$H=\langle A^2,B^2\rangle$, and tests all
$\lambda\in(\mathbb Z/\exp G)^\times$, $f\in H'$.  The results are
\[
\begin{array}{c|r|r|r|r}
 n&|G|&\text{candidate pairs}&\text{braid solutions}&
 \text{quotient automorphisms}\\\hline
 2&6&2&2&2\\3&24&32&32&24\\4&48&4&4&4\\
 5&120&1920&960&800\\6&144&32&32&24\\
 7&336&16128&5376&4704\\8&384&16&16&16\\
 9&648&2592&720&648\\10&720&1920&960&800
\end{array}
\]
The next modulus has order $1320$.  Every successful quotient map is
verified on all Cayley edges, including well-definedness on different
words for the same element.  All $2n$ choices obtained from
$\lambda=\pm1,f=1$ and the conjugation $A^{-8k}$, $k\bmod n$, pass
for each row.  The excess in the fourth column proves that a braid
relation test alone does not even certify descent to the chosen
quotient.  These candidates are not labelled GT-shadows: the
arity-four pentagon and compatible quotient data have not been
supplied for that table.  A classification over \emph{all} finite
quotients of $B_3$ of order at most $1000$, or of the non-charming
fake shadows, is not performed here.  This is the precise remaining
scope if the numerical request is read in that fully exhaustive sense.

\subsection{Covering germs, tangential fibres, and the proposed limiting topos}

\begin{statement}{E.10}{Proved}{The covering germ of a puncture and its tangent fibre functors}
Let $X$ be an oriented topological surface and $p\in X$.  For a
cofinal system of chart disks $D_n$ shrinking to $p$, put
$D_n^*=D_n\setminus\{p\}$.  With morphisms defined after a common
restriction, the filtered category of covering germs is
\[
 \operatorname*{2\!\varinjlim}_n\operatorname{Cov}(D_n^*)
       \simeq\mathbb Z\text{-}\operatorname{Set}.
\]
Its finite-cover subcategory is equivalent to finite continuous
$\widehat{\mathbb Z}$-sets.  A topological orientation reversal
acts on the local monodromy by $k\mapsto-k$.
For a $C^1$ surface, a nonzero tangent ray defines a fibre functor
by an outgoing $C^1$ arc.  The groupoid of these fibre functors
and their path transports is the fundamental groupoid of the circle
of tangent rays: each transport set is a $\mathbb Z$-torsor.
On finite covers the torsors of all natural isomorphisms are their
profinite completions; geometric path transports form the dense
integer subtorsors.
The same local statement holds near a smooth complex divisor in
coordinates $\Delta^{d-1}\times\Delta^*$, and near a normal-crossing
point with $r$ branches it has $\mathbb Z^r$ in place of $\mathbb Z$.
\end{statement}
\begin{proof}
In polar coordinates the universal cover of a punctured chart disk
is $(0,\epsilon)\times\mathbb R$, with deck transformation
$(s,t)\mapsto(s,t+1)$.  An arbitrary $\mathbb Z$-set $E$ gives the
cover $((0,\epsilon)\times\mathbb R)\times_{\mathbb Z}E$.
Lifting paths proves that every covering has this form and that
morphisms are exactly equivariant maps.  Every inclusion of one
concentric punctured disk into another induces the identity on the
oriented fundamental group.  Its restriction functor on covering
categories is therefore an equivalence.  This proves the filtered
statement, for arbitrary fibre sets as well as for finite ones,
without exchanging infinitely many local triviality conditions.
Finite $\mathbb Z$-actions uniquely extend continuously to
$\widehat{\mathbb Z}$.  A local homeomorphism of surface germs has
degree $+1$ or $-1$ on a small linking circle, giving the sign assertion.

For a nonzero tangent $v$, choose a chart with $p=0$ and an arc
$\gamma(t)=tv+o(t)$.  Pullback to the interval gives a trivial
covering; its connected lift components define the fibre functor.
On a sufficiently short interval,
$|\gamma(t)-tv|<|v|t/2$.  The straight homotopy from $\gamma(t)$
to $tv$ stays outside $0$.  The same estimate for two arcs and
for their homotopies proves independence of the arc with that
first-order direction, up to the specified coherent transport.
A path of nonzero tangent rays gives transport by a small circular
arc, and two such paths give the same transport exactly when their
winding difference is zero.  One can test the latter assertion on
the universal $\mathbb Z$-cover.  This proves the assertion about
torsors and groupoids.  If one tests only finite covers, every
finite quotient of $\mathbb Z$ is available, so the automorphism
group of the fibre functor is $\widehat{\mathbb Z}$.
Finally the remaining coordinate factors in the divisor models
are contractible, while each punctured coordinate contributes one
independent winding number.  This proves the asserted local forms
of Lead~4 and of scans 242--272; it asserts no global divisor theorem
without corresponding neighbourhood hypotheses.
\end{proof}

\begin{statement}{E.11}{Refuted}{The full small-topos inverse limit loses the punctured germ}
For $D_n^*=\{z\in\mathbb C:0<|z|<1/n\}$, the literal inverse
$2$-limit of the full small etale topoi
$\operatorname{Sh}(D_n^*)$, with the open-inclusion geometric
morphisms, is the degenerate topos.  Its covering category therefore
cannot be the category in E.10.  In particular the parenthetical
candidate at C4 scans 242--243, fols.~768'--769, is not a correct
identification for this basic residual surface germ when ``small
etale topos'' has its usual full-sheaf meaning.

A correct topos encoding its covering germs is the classifying
topos $B\mathbb Z=\mathbb Z\text{-}\operatorname{Set}$.  It has the
small site consisting of the one-object group $\mathbb Z$ with the
trivial topology.  Its finite covering data are the finite continuous
$\widehat{\mathbb Z}$-sets.  One can retain the original neighbourhood
information as the pro-topos $(\operatorname{Sh}(D_n^*))_n$, without
replacing it by this degenerate inverse limit.
\end{statement}
\begin{proof}
Let a topos $\mathcal E$ have a compatible cone of geometric
morphisms $f_n:\mathcal E\to\operatorname{Sh}(D_n^*)$.
For each $n$, viewed as an open subobject of $D_1^*$, the subset
$D_n^*$ pulls back along $f_1$ to the terminal object: $f_1$ factors
through its open subtopos.  Put
$U_k=\{z:1/k<|z|<1\}\subset D_1^*$, $k\geq2$.
It is disjoint from $D_k^*$, so finite-limit preservation gives
$f_1^*(U_k)=0$.  But $\bigcup_{k\geq2}U_k=D_1^*$.
Inverse-image functors preserve all unions of subobjects (they
preserve coproducts, finite limits and their images), so
\[
 1_{\mathcal E}=f_1^*(D_1^*)=
       \bigvee_{k\geq2}f_1^*(U_k)=0_{\mathcal E}.
\]
Thus every such cone has degenerate source.  The degenerate topos
itself has a unique compatible cone, proving the claimed universal
property of the inverse $2$-limit.
Equivalently, present its locale by generators $[U,n]$ for all
opens $U\subset D_n^*$, with the stagewise frame relations and
$[U,n]=[U\cap D_m^*,m]$ for $m\geq n$.  This is a small presentation
of the colimit frame.  Each $U_k$ becomes zero and their join is
one, so the frame is trivial.  This also constructs the requested
literal limit rather than merely observing that the set-theoretic
intersection of the spaces is empty.

E.10 supplies the nondegenerate replacement $B\mathbb Z$.  In this
example every covering and every morphism descends to an individual
disk because each restriction functor is already an equivalence.
This is a direct proof of finite descent for the covering-germ
system.  It is not a compactness argument for descent from the
full-sheaf limit, which the first half of the proof rules out.
The manuscript explicitly qualifies its expectation by ``good cases'';
our refutation concerns applying its indicated full-topos construction
to the punctured-disk case at hand.  Neither this limit nor its
covering-topos replacement is a tower of Lie approximations.
\end{proof}

\begin{statement}{E.12}{Known}{The valid differential/topological comparison and its low-type exceptions}
For a compact smooth oriented surface $X$ with a finite labelled
set $I$, inclusion of smooth diffeomorphisms into homeomorphisms
induces the same mapping class group.  If $2g-2+|I|>0$, their
identity components are weakly contractible; the diffeomorphism
component is contractible.  The lifted action on punctured-neighbourhood
coverings gives the same component group as the lifted differential
action on tangent rays.  It fits into
\[
 1\to\mathbb Z^I\to ST_{g,I}\to\operatorname{Mod}^{\pm}(X,I)\to1,
\]
with orientation signs and, when labels may be permuted, the induced
permutation action on the kernel.

Without the hyperbolicity hypothesis, this kernel assertion is false.
For the sphere with one marked point the positive lifted component
group is trivial.  For the sphere with two individually fixed marked
points it is $\mathbb Z$, with kernel
$\mathbb Z^2/\langle(1,-1)\rangle$, not $\mathbb Z^2$.
For $(1,0)$ there is no tangent kernel and the group is
$\mathrm{GL}_2(\mathbb Z)$; for $(0,0)$ it is $C_2$.
These qualifications apply to the preliminary assertion on scan 227;
the hyperbolic restriction is made explicit later on scan 277.
\end{statement}
\begin{proof}[Theorems used and comparison argument]
For diffeomorphisms the base cases are
\citet[Theorem~1(a)--(c)]{EarleEells1967}: a sphere with three fixed
points, a torus with one fixed point, and a closed surface of genus
at least two.  The corresponding homotopy statements for
homeomorphisms are proved and reviewed by
\citet[Theorems~1.5.1--1.5.2 and \S1.7]{Hamstrom1974}.
For the marked torus and sphere cases, evaluate respectively on one
point and on an ordered triple.  Translations and Mobius transformations
supply the sections and the noncontractible factors, so removing
these factors leaves a weakly contractible stabilizer.  Adding more
marked points uses the evaluation fibration with base $X\setminus I$.
Once in one of the three base cases this surface is aspherical.
The homotopy exact sequence shows that the identity component of
the new stabilizer remains weakly contractible.  In the smooth category
these stabilizers are metrizable Frechet manifolds (use local vector
fields vanishing at the finitely many markings).  They are ANRs by
\citet[Theorem~5]{Palais1966}, and hence weak contractibility implies
contractibility by his Theorem~15 and its corollary.
For surjectivity on component groups, the precise smoothing input is
\citet[Theorem~B and the Handle Smoothing Theorem]{Hatcher2022Torus}.
One first straightens near the finitely many markings, retaining
their centres, and applies its relative version on the remaining
surface.  For injectivity the cut-graph proof of the surface isotopy
theorem applies in both categories: homotopic embedded cut arcs are
made disjoint by successive innermost-bigon removals and then moved
onto one another; after fixing the graph one uses the disk isotopy
theorem on the complementary disks.  In the smooth category these
moves are supported in smooth disks and strips and extend by smooth
ambient isotopies.  The topological arc and relative approximation
steps are \citet[Theorem~3.1 and Appendix, Theorems~A1--A4]{Epstein1966}.
Use a cut graph containing $I$ among its vertices to retain all
markings throughout.  This proves the component comparison rather
than assuming it from contractibility of the identity components.
No assertion equating individual homeomorphisms and diffeomorphisms
is needed.

Here is the compatibility relevant to the manuscript.  If $h$ is a
$C^1$ germ fixing $0$ with invertible derivative, then
\[
 h_s(v)=s^{-1}h(sv)\quad(s>0),\qquad h_0(v)=D_0h(v)
\]
is an isotopy of punctured germs, continuous in $C^1$ on a sufficiently
small common neighbourhood.  Its inverse germs satisfy the same
formula.  It identifies pullback on covering germs with that of the
linear derivative, and E.10 identifies the latter with its action
on tangent rays.  Composition is compatible because
$(h\circ k)_s=h_s\circ k_s$.
For homeomorphisms the punctured covering action is defined without
a derivative and a lift has the same integer ambiguity.  In a family
near the identity this is a covering of the transformation group:
choose a small punctured chart and follow a fixed outgoing arc;
path lifting gives local branches, whose differences are integer
deck transformations.  Its restriction to diffeomorphisms is the
derivative covering just described.  The exact homotopy sequence is
\[
 \pi_1(\operatorname{Homeo}_0(X,I))\longrightarrow\mathbb Z^I
 \longrightarrow\pi_0(\widetilde{\operatorname{Homeo}}(X,I))
 \longrightarrow\pi_0(\operatorname{Homeo}(X,I))\to1,
\]
and likewise in the smooth category.  Weak contractibility kills
the first term.  The mapping class comparison and the identity map
on $\mathbb Z^I$ therefore identify the two lifted extensions.

For the one-point sphere the identity component retracts to the
circle of rotations fixing the point; one full rotation winds its
tangent circle once.  The first arrow in the last display is then
$\mathbb Z\xrightarrow{1}\mathbb Z$, killing the proposed kernel.
For two labelled points the identity component again has the rotation
circle.  In local oriented complex coordinates $z$ and $w=1/z$ its
rotation winds by $(1,-1)$, giving the stated quotient.  When $I$ is
empty there is no lift parameter at all, yielding the remaining
mapping class groups.  These computations explain exactly which
step of the hyperbolic proof fails in the exceptional cases.
\end{proof}

\begin{statement}{E.13}{Refuted}{The unrestricted germ assertions require genuine qualifications}
The inclusion of differentiable surface germs into homeomorphism
germs is not an equality, and topological germs need not induce a
map on tangent rays.  Also the assertion on scan 245 that the
ordinary germ at a point is always simply connected is false for
arbitrary topological spaces if the category includes all infinite
coverings and all their morphisms.  It is correct for the locally
connected surface germs used in E.10, and the arbitrary-point
assertion is correct for finite covering germs.
\end{statement}
\begin{proof}
The radial homeomorphism $h(z)=z|z|$ has derivative zero at the
origin and is not a local diffeomorphism.  The homeomorphism
$h(re^{i\theta})=re^{i(\theta+\log r)}$, extended by $h(0)=0$,
has an inverse of the same form and is continuous at $0$, but
its outgoing rays have no limiting tangent direction.  Both still
act on the covering germ through the orientation-preserving
monodromy of E.10.  Thus it is the covering and isotopy comparison,
not literal equality of these germ groups, which is valid.

For the general-point assertion let
$X=\{0\}\cup\{1/n:n\geq1\}$ and take its germ at $0$.
Consider the constant coverings with fibres $E=\mathbb N$ and
$F=\{0,1\}$.  Define a map over $X$ by sending sheet $j$ to $1$
only over $1/j$, and to $0$ everywhere else.  On each sheet this
is continuous: the singleton $\{1/j\}$ is clopen in $X$.
At $0$ all sheets map to $0$.  Nevertheless on every neighbourhood
of $0$ some sheet maps to $1$ somewhere, so this morphism germ
is not the germ of any single set map $E\to F$.  The constant
covering functor from sets is not full.  Every covering is locally
trivial at $0$, so the failure is in morphisms, exactly the missing
uniformity across infinitely many sheets.  For finite fibres
one can intersect the finitely many neighbourhoods needed for
constancy of a morphism; local triviality of objects then proves
that the finite covering germ category is finite sets.  Surface
neighbourhoods have a cofinal system of connected disks, so the
same uniformity issue does not arise even for infinite fibres there.
\end{proof}
